\documentclass[11pt]{article}
\usepackage{setspace}
\usepackage{geometry}
\geometry{margin=1in}
\setstretch{1.15}

\begin{document}

% --- Applicant Information ---
\noindent
\textbf{Zhongli Wang} \\
Phone: +49 160 9541 6015 \\
Email: nkwangzhongli@gmail.com \\
Website: https://sites.google.com/view/zhongli-wang \\
Bielefeld, Germany \\

\vspace{0.5cm}

% --- Recipient Information ---
\noindent
\textbf{Dr.\ Aleksandrina Ralcheva and Hiring Committee} \\
Innovation for Delft Engineering Education (IDEE) \\
TU Delft \\
Delft, Netherlands \\
\vspace{0.5cm}

% --- Subject Line ---
\noindent
\textbf{Application for Postdoctoral Researcher — Teaching and Learning in the Age of AI} \\

% \vspace{0.3cm}

\noindent Dear Dr.\ Ralcheva and Members of the Hiring Committee,\\[0.2cm]

I am writing to apply for the Postdoctoral Researcher position on “Teaching and Learning in the Age of AI” within the IDEE initiative at TU Delft.  
I am currently completing my Ph.D.\ in Economics at Bielefeld University and Università Cattolica del Sacro Cuore through the Marie Skłodowska-Curie EPOC program.  
My research focuses on human–AI interaction, computational modeling of learning processes, and the development of AI-enabled workflows that support analytical reasoning and knowledge acquisition.  
Across my projects, I study how individuals form judgments, adapt strategies, and learn in AI-rich environments—using methods that combine machine learning, behavioral modeling, experimental design, and analysis of human–AI collaboration.  

During my PhD, which I began in 2021, I increasingly shifted my research toward understanding how individuals learn, make decisions, and adapt within AI-rich digital environments.  
To study these processes, I developed computational and data-driven models that examine how users interact with algorithmic systems, how information structures shape learning trajectories, and how behavioral patterns emerge in technology-mediated contexts.  
Across these projects, I worked extensively with large datasets, Python-based simulation environments, and mixed-method analytical pipelines—skills directly transferable to studying student–AI interaction and meaningful learning processes.  
These research designs required integrating theory with empirical evidence, validating behavioral assumptions, and translating model outputs into insights relevant for instructional design and learning outcomes.  
All three chapters of my dissertation are single-authored, demonstrating my ability to independently design, execute, and communicate complex research projects from conceptualization to implementation.

Beyond my dissertation work, I have coauthored two projects that demonstrate my ability to apply machine learning and AI-driven analytics to understand and design complex human–AI systems. 
One project, “Polytope Fraud Theory,” uses machine learning and statistical methods to detect patterns and anomalies in large-scale datasets, illustrating my experience with high-dimensional data analysis, feature engineering, and model evaluation—competencies directly transferable to studying learning traces and AI-mediated student behaviors. 
The second project develops an agent-based workflow for integrating large-language-model agents into empirical research, creating automated pipelines for text analysis, reasoning tasks, and human–AI collaboration. 
This work familiarized me with the opportunities and limitations of LLM-assisted analysis and highlighted how system design shapes user behavior, cognitive load, and learning outcomes. 
Together, these experiences strengthened my ability to build scalable analytical tools, design reproducible AI workflows, and translate technical outputs into insights that support evidence-based improvements in teaching and learning.

I am highly motivated to join TU Delft because of its leading role in rethinking education in the age of AI and its strong culture of interdisciplinary collaboration. 
I am particularly drawn to the opportunity to work closely with educators, learning scientists, engineers, and policymakers to understand how AI reshapes teaching, learning processes, and assessment design. 
My background in AI-enabled analysis and computational modeling has prepared me to approach open-ended educational questions with methodological rigor, careful attention to human–AI interaction, and a commitment to generating insights that are both theoretically grounded and practically actionable. 

At TU Delft, I hope to contribute my computational, analytical, and behavioral modeling expertise to advancing research on how AI transforms teaching and learning. 
My areas of particular interest include:
\begin{itemize}
    \item Analyzing patterns of student–AI interaction to distinguish deep learning from surface learning and to identify conditions that promote meaningful engagement;
    \item Developing empirical and computational frameworks to evaluate how AI tools influence learning objectives, teaching activities, and assessment strategies within constructive alignment;
    \item Applying qualitative and quantitative methods—including surveys, interviews, experiments, and learning analytics—to evaluate AI-supported learning environments and inform evidence-based educational design.
\end{itemize}

% I am especially motivated by Apple’s commitment to analytical rigor, innovation, and cross-functional collaboration. 
% I would welcome the opportunity to contribute to Apple’s strategic decision-making by transforming complex quantitative analysis into clear, actionable insights that support growth, performance optimization, and long-term planning.

Thank you so much for considering my application! I would welcome the opportunity to discuss how my analytical, computational, and research communication skills can contribute to the IDEE initiative at TU Delft, and how I can support the development of evidence-based frameworks, guidelines, and research on meaningful learning in the age of AI.

\vspace{5mm}

\noindent Sincerely,\\[15mm]
\noindent Zhongli Wang



\end{document}